% BHU PYQ -- Manually transcribed & error-corrected
% Paper: BCF-215 : Corporate Laws
% Exam: B.Com. (Hons.) F.M.M. Semester III Examination, 2016-17

\documentclass[11pt,a4paper]{article}
\usepackage[a4paper,top=2.5cm,bottom=2.5cm,left=2.8cm,right=2.8cm,headheight=14pt]{geometry}
\usepackage{amsmath,amssymb}
\usepackage[T1]{fontenc}
\usepackage[utf8]{inputenc}
\usepackage{lmodern,microtype,fancyhdr,enumitem,hyperref,lastpage,parskip}

\hypersetup{
    colorlinks=true,
    urlcolor=black,
    linkcolor=black,
    pdftitle={BCF-215: Corporate Laws -- BHU B.Com. (Hons.) FMM Sem III 2016-17}
}

\pagestyle{fancy}
\fancyhf{}
\renewcommand{\headrulewidth}{0.4pt}
\renewcommand{\footrulewidth}{0.4pt}
\fancyhead[L]{\small   \textit{Banaras Hindu University}}
\fancyhead[R]{\small   \textit{BCF-215: Corporate Laws}}
\fancyfoot[C]{\small Page  \thepage\ of \pageref{LastPage}}


\newlist{parts}{enumerate}{2}
\setlist[parts,1]{label=(\alph*),leftmargin=*,itemsep=5pt,topsep=3pt}
\setlist[parts,2]{label=(\roman*),leftmargin=*,itemsep=3pt,topsep=2pt}

\newcommand{\pts}[1]{\hfill   \textit{[#1]}}


\begin{document}


\begin{center}
  {\large\bfseries BANARAS HINDU UNIVERSITY}\\[2pt]
  {\normalsize B.Com. (Hons.) F.M.M. --- Semester III Examination, 2016-17}\\[6pt]
  \rule{0.8\linewidth}{0.5pt}\\[6pt]
  {\large\bfseries Paper No. BCF-215: Corporate Laws}\\[4pt]
  {\normalsize Time: 3 Hours\hfill Maximum Marks: 70}
\end{center}

\rule{\linewidth}{0.4pt}

\smallskip



\noindent   \textit{Instructions: Answer all questions. All questions carry equal marks (14 marks each).}

\bigskip




\noindent   \textbf{Question 1.}
``A company is an artificial person created by law with a perpetual succession and a common seal.'' Explain this statement and also describe the conditions under which the corporate veil of a company can be lifted.\pts{14}

\centerline{   \textbf{OR}}

What is a private company? Describe the statutory rights and exemptions available to a private company.\pts{14}

\medskip\rule{\linewidth}{0.2pt}




\noindent   \textbf{Question 2.}
What is meant by promotion of a company? What formalities must be complied with in the formation of a company under the Indian Companies Act?\pts{14}

\centerline{   \textbf{OR}}

Describe the procedure relating to public company registration. Can the registration certificate be cancelled under any circumstance?\pts{14}

\medskip\rule{\linewidth}{0.2pt}




\noindent   \textbf{Question 3.}
Distinguish between the Memorandum and the Articles of Association.\pts{14}

\centerline{   \textbf{OR}}

What is Articles of Association? What are its contents?\pts{14}

\medskip\rule{\linewidth}{0.2pt}




\noindent   \textbf{Question 4.}
What is meant by director? State the conditions under which a person is disqualified from being appointed as a director.\pts{14}

\centerline{   \textbf{OR}}

What is meant by company secretary? How is a company secretary appointed?\pts{14}

\medskip\rule{\linewidth}{0.2pt}




\noindent   \textbf{Question 5.}
Describe the duties and liabilities of a company secretary.\pts{14}

\centerline{   \textbf{OR}}

State and explain the different kinds of meetings of a company.\pts{14}

\vfill
\rule{\linewidth}{0.4pt}

\begin{center}\small
\href{https://bhu-exam.vercel.app/}{Click here to visit our website}\\[4pt]
\href{https://me-aryan.vercel.app/}{Click here to visit our contributor}\\[4pt]
\href{https://quashan-me.vercel.app/}{Click here to visit our contributor}
\end{center}

\end{document}
